% Higher UK long-term interest rates --- seminar slides
% Calibrated small-open-economy NK twin in uk-long-rate/
%
% Compile from this directory:
%   pdflatex uk_long_rate_seminar.tex
%   tectonic uk_long_rate_seminar.tex
%
% Figures are the slide-ready PDFs in ../output/slides/
% (src/plot_slides.py). Change \irfdir only.
\documentclass[11pt,aspectratio=169]{beamer}

\usetheme{Madrid}
\usecolortheme{default}
\setbeamertemplate{navigation symbols}{}
\setbeamertemplate{itemize items}[circle]
\setbeamerfont{footnote}{size=\tiny}

\usepackage[T1]{fontenc}
\usepackage[utf8]{inputenc}
\usepackage{amsmath,amssymb}
\usepackage{booktabs}
\usepackage{graphicx}
\usepackage{array}

\newcommand{\irfdir}{../output/slides/}
\graphicspath{{\irfdir}}

\newcommand{\Calibrated}{%
  \par\vspace{0.2em}%
  {\footnotesize\textbf{Calibrated model.}
  Bayesian estimation has not been run.
  $\kappa$ and the housing parameters are provisional.}%
}

\title[UK long rates]{Higher UK Long-Term Interest Rates}
\subtitle{\texorpdfstring{Small-open-economy NK with a HANK-lite mortgage block\\
{\normalsize Seminar notes on the \texttt{uk-long-rate} MVP}}{UK long-rate seminar notes}}
\author[Prototype slides]{\texorpdfstring{Open-source notes on the signed sketch and the calibrated twin\\
{\small not a Bayesian estimate}}{Open-source prototype slides (calibrated, not estimated)}}
\institute[MacroModellerBot]{\texttt{uk-long-rate/} in MacroModellerBot}
\date{September 2026}

\begin{document}

%----------------------------------------------------------------------
\begin{frame}[plain]
  \titlepage
  \vspace{-1.1em}
  \begin{center}
    {\footnotesize
    Sources: signed channels review; model sketch v3 (21~Sep 2026);\\
    Dynare twin \texttt{dynare/uk\_long\_rate.mod} and the Python QZ paths in \texttt{output/}.}
  \end{center}
\end{frame}

%----------------------------------------------------------------------
\begin{frame}{Outline}
  \tableofcontents
  \vspace{0.6em}
  {\footnotesize Channels from the signed review first, then the calibrated
  twin. Every results path below is that twin.
  Bayesian estimation has not been run.}
\end{frame}

%======================================================================
\section{The question}
%----------------------------------------------------------------------
\begin{frame}{What we are asking}
  \begin{block}{Question (signed sketch)}
    How does an exogenous rise in UK long-term nominal rates transmit
    to GDP, CPI, house prices, sterling, consumption, investment and
    the public interest bill, with the short policy-rate path held
    fixed or allowed to move?
  \end{block}
  \vspace{0.2em}
  Three shocks, kept separate. Illustrative size in the model:
  a $+100$~bp move in the model 10-year rate.
  \begin{itemize}
    \item \textbf{(a) sterilised term premium.}
          The term premium lifts the long rate. The expected Bank Rate
          path is held fixed.
    \item \textbf{(b) path news.}
          News about future short rates. Bank Rate is not pegged.
    \item \textbf{(c) sovereign.}
          A fiscal / sovereign gilt risk premium.
          \textbf{Not yet modelled.}
  \end{itemize}
\end{frame}

%----------------------------------------------------------------------
\begin{frame}{Why ``the 10-year yield'' is not one shock}
  \begin{itemize}
    \item A long yield is an expected path of short rates plus a term
          premium (duration, inflation, liquidity).
    \item A path-driven rise often co-moves with stronger expected
          activity or inflation. A term-premium rise --- gilt supply,
          QE/QT, fiscal credibility, risk-off --- tightens more like a
          cost of capital
          (Kaminska, Mumtaz and \v{S}ustek, 2021;
          Greenwood, Hanson, Stein and Sunderam, 2023).
    \item High-frequency surprises are not the same object as the
          gradual 2021--23 repricing, which was partly anticipated and
          front-loaded into mortgages (Burr and Willems, 2024).
    \item For the UK, long rates price fixed-rate mortgages (via swaps),
          the public interest bill, and an open-economy portfolio balance
          into sterling.
  \end{itemize}
  \vspace{0.35em}
  {\footnotesize Do not treat ``10-year yield up 100~bp'' as interchangeable
  across \textbf{(a) sterilised term premium}, \textbf{(b) path news}
  and \textbf{(c) sovereign}.}
\end{frame}

%======================================================================
\section{Literature}
%----------------------------------------------------------------------
\begin{frame}{Path versus term premium}
  \begin{itemize}
    \item \textbf{Comparable UK macro IRFs} for a pure term-premium shock
          versus a Bank Rate / path shock are sparse
          (signed review, \S1.1).
    \item UK \emph{path} peaks are the high-frequency Bank Rate / MPC
          studies. UK \emph{term-premium} evidence is much better for
          \emph{yields} (Joyce and Lengyel, SWP~1,097) than for full
          macro IRFs.
    \item Kaminska, Mumtaz and \v{S}ustek (SWP~914, 2021; also \emph{JME})
          is a \textbf{US} FOMC decomposition. Useful for the
          \emph{shape} of path versus premia responses. Instruments are
          normalised differently. Not a UK calibration target.
    \item The review's modelling implication, which this MVP follows:
          separate \textbf{(a) sterilised term premium},
          \textbf{(b) path news} and \textbf{(c) sovereign}.
          Empirically, sterilise (a) by residualising long-end surprises
          on Bank Rate / short OIS.
  \end{itemize}
\end{frame}

%----------------------------------------------------------------------
\begin{frame}{Magnitudes, with the shock labelled}
  {\scriptsize
  \setlength{\tabcolsep}{3.5pt}
  \renewcommand{\arraystretch}{1.12}
  \begin{tabular}{@{}>{\raggedright\arraybackslash}p{2.55cm}>{\raggedright\arraybackslash}p{3.15cm}>{\raggedright\arraybackslash}p{6.35cm}@{}}
    \toprule
    Object & Source & Peak / key magnitude \\
    \midrule
    Bank Rate / MPC surprise
      & Cesa-Bianchi, Thwaites and Vicondoa (2020); ranges in Burr and Willems (2024). UK.
      & GDP roughly $-0.5\%$ to $-1.4\%$; prices $-0.3\%$ to $-1.0\%$; unemployment $\sim +0.5$~pp in some estimates; houses $\sim -4\%$; sterling ERI $\sim +2\%$ to $+7\%$. Per $\approx 1$~pp. \textbf{Path}, not pure term premium. \\
    Housing share of consumption
      & Albuquerque, Lazarowicz and Lenney, SWP~\textbf{1,115} (2025). UK.
      & Housing channel $\approx$ \textbf{one half} of peak transmission to household consumption, among mortgagors. \textbf{Path / rate} shock. \\
    Gilt supply / QT-sized surprise, 10-year yield
      & Joyce and Lengyel, SWP~\textbf{1,097} (2024). UK.
      & $\sim 20$~bp in low stress; $\sim 120$~bp in high stress. Yield impact, not a macro IRF. \textbf{Term premium / gilt supply}. \\
    Action, expected path, uncertainty
      & Kaminska, Mumtaz and \v{S}ustek, SWP~\textbf{914} (2021). \textbf{US}.
      & Action: NK contraction. Path: often expansionary IP. Premia: yields and mortgage rates up, housing down. Shapes only. \\
    \bottomrule
  \end{tabular}
  }
  \vspace{0.25em}
  {\scriptsize Condensed from the signed review, \S1.1.
  Passive QT runoff $\approx$ active sales, pound for pound, in SWP~1,097.}
\end{frame}

%----------------------------------------------------------------------
\begin{frame}{Mortgage cash flow: 2-year and 5-year fixes}
  \begin{itemize}
    \item New UK fixes are priced off sterling swaps, which co-move with
          gilts. A gilt spike can reprice the \emph{flow} between MPC
          meetings (Burr and Willems, 2024).
    \item The \emph{stock} is slow. The share of outstanding mortgages
          fixed for five years rose from around \textbf{20\% in 2016 to
          over 50\% in 2022} (Burr and Willems, 2024, fn.~18).
    \item Over 90\% of \emph{new} mortgages have been fixed-rate since
          around 2013, with five-year products replacing two-year fixes
          (Burr and Willems, 2024; MLAR flow evidence).
    \item Cash-flow peaks should lag the peak in flow mortgage rates,
          over a \textbf{2--5 year} reset ladder. Anticipation of a
          dearer refinance can move some spending earlier
          (Anjum and Herler, 2023, \emph{cited via} Burr and Willems).
    \item Effective stock rates rose more slowly in 2021--23 than in
          past cycles because of that fixed share (IMF, 2024).
  \end{itemize}
  \vspace{0.15em}
  {\scriptsize SWP~\textbf{1,158} (Rajan, Rodriguez-Tous and Salgado-Moreno, 2025)
  studies how term \emph{spreads} shift the supply of short versus long
  fixes. It is not a source for stock reset shares.}
\end{frame}

%----------------------------------------------------------------------
\begin{frame}{Housing collateral, alongside the cash-flow channel}
  \begin{itemize}
    \item Higher long rates raise the user cost of housing, cut the
          present value of housing services, and tighten LTVs when
          prices fall (Mishkin, 2007; IMF WEO, April 2024, Ch.~2).
    \item BoE staff: a 1~pp \textbf{Bank Rate} hike lowers UK house
          prices by roughly 4\% on average
          (Burr and Willems, 2024). That is a \textbf{path} magnitude,
          not a term-premium IRF.
    \item SWP~\textbf{1,115}: the housing channel is about half of peak
          transmission to household consumption, concentrated among
          mortgagors. In their high-frequency evidence, house prices
          trough around one year; nominal rents are stable for one to
          two years, then decline.
    \item Housing wealth effects on consumption are typically at least
          as large as equity wealth effects
          (Case, Quigley and Shiller, 2001). UK MPCs out of wealth
          are cited at around 5p per \pounds1
          (Slacalek, 2009, \emph{cited via} Burr and Willems).
    \item Direct intertemporal substitution looks small next to
          cash flow, income and wealth
          (Havranek et al., 2015; Best et al., 2020, for the UK EIS).
  \end{itemize}
\end{frame}

%----------------------------------------------------------------------
\begin{frame}{Investment, and corporate credit versus gilt premia}
  \begin{itemize}
    \item User cost rises with the real financing rate
          (Hall and Jorgenson, 1967). Long-lived capital is the usual
          margin; timing can also respond to short rates
          (McKay and Wieland, \emph{cited via} Burr and Willems, 2024).
    \item Gilchrist and Zakraj\v{s}ek (2007): a 1~pp rise in user cost
          cuts the investment rate by about 50--75~bp in the short run
          (roughly 30--50~bp from a 1~pp rise in the real long rate),
          with a long-run capital-stock elasticity near one.
    \item UK firm evidence in the hiking cycle: higher financing costs
          plus weaker demand (Burr and Willems, 2024; Shah et al., 2024,
          \emph{cited via} them). Internal cash buffers dampen the
          response (Joseph et al., 2019, \emph{cited via} them).
          Corporate-bond pass-through settled near $\sim 0.7$
          (IMF UK Selected Issues, 2024).
    \item A gilt term premium is not an excess bond premium.
          Gilchrist and Zakraj\v{s}ek (2012) split US spreads into
          expected default and an EBP; Gertler and Karadi (2015)
          stress credit costs as well as the funds rate.
  \end{itemize}
  \vspace{0.2em}
  {\scriptsize The review allows a thin spread block, or an explicit
  deferral. This MVP defers IG/HY credit (and CRE). Transmission runs
  through gilt / swap user costs.}
\end{frame}

%----------------------------------------------------------------------
\begin{frame}{The public interest bill}
  \begin{itemize}
    \item Higher long yields raise debt service on new and refinanced
          gilts, and bond supply absorbed by risk-averse investors can
          crowd out private investment
          (Greenwood and Vayanos, 2014).
    \item Joyce and Lengyel (SWP~\textbf{1,097}, 2024): issuance
          surprises move the curve through duration-risk and local-supply
          channels. A fully unexpected QT-sized programme:
          about \textbf{20~bp} at the 10-year in low stress, nearly
          \textbf{120~bp} in high stress. Passive runoff and active
          sales are broadly similar, pound for pound.
    \item IMF (2024): gilt yields were more responsive to policy-rate
          increases than in past UK cycles, consistent with elevated
          debt and post-mini-Budget risk perceptions.
    \item With elevated public debt, a term-premium shock has a
          first-order fiscal leg, and that leg can feed back into
          confidence and gilt-market functioning.
  \end{itemize}
  \vspace{0.3em}
  {\footnotesize Modelling implication from the sketch: keep an
  index-linked share. Do not set $s^{IL}=0$.}
\end{frame}

%----------------------------------------------------------------------
\begin{frame}{Sterling depends on the shock; LDI is not linear}
  \begin{columns}[T]
    \begin{column}{0.48\textwidth}
      \textbf{FX sign (review \S2.5)}
      \begin{itemize}
        \item A 1~pp \textbf{Bank Rate} surprise:
              sterling ERI about $+2\%$ to $+7\%$,
              peaking within about a quarter
              (Burr and Willems, 2024; Forbes et al.,
              \emph{cited via} them).
        \item \textbf{(a) sterilised term premium:}
              a relative-rate reading.
              Sterling tends to \textbf{appreciate}.
        \item \textbf{(c) sovereign} (review only):
              gilt--FX risk, so the same long-rate
              rise co-moves with a
              \textbf{depreciation}.
              Not yet modelled here.
      \end{itemize}
      {\scriptsize Do not hard-wire ``higher long rates
      $\Rightarrow$ a stronger pound''
      (Greenwood et al., 2023;
      Gourinchas and Rey, 2007).}
    \end{column}
    \begin{column}{0.50\textwidth}
      \begin{alertblock}{LDI: non-linear stress, off in the MVP}
        {\scriptsize
        September 2022: leveraged LDI faced NAV and collateral
        calls after the 23~September mini-Budget; forced gilt sales
        amplified the sell-off
        (BoE Financial Policy Summary, 30~Sep 2022;
        Financial Stability Report, December 2022;
        Pinter, SWP~\textbf{1,019}, 2023).
        Joyce and Lengyel's $\sim 20$ versus $\sim 120$~bp gap is
        the same stress non-linearity.
        \textbf{Leave it off} for a sterilised, calm-market term
        premium. Switch it on only under gilt stress, leveraged
        non-bank duration, or a large sudden credibility shock.}
      \end{alertblock}
    \end{column}
  \end{columns}
\end{frame}

%======================================================================
\section{The model}
%----------------------------------------------------------------------
\begin{frame}{SOE-NK with HANK-lite savers and mortgagors}
  \begin{itemize}
    \item Small open economy. Hybrid Phillips curve, Taylor rule,
          $q$-theory investment, UIP.
    \item Two households only: \textbf{savers} and \textbf{mortgagors}.
          Not a full idiosyncratic-income distribution.
    \item Mortgage bridge from the yield curve. One remortgage hazard
          updates both the stock rate and the LTV debt stock.
          No second cohort, no buy-to-let.
    \item Sticky housing user-cost curve, so house prices are not an
          impact jump. Fiscal interest bill with an index-linked share
          and a duration refix. Resource constraint with net exports.
    \item Linear system in \texttt{dynare/uk\_long\_rate.mod}:
          19 predetermined variables, 6 forward-looking variables in
          the declaration, 11 static definitions.
          Blanchard--Kahn count on this calibration: 36 stable roots,
          6 finite unstable (\texttt{output/success\_checks.txt}).
    \item The paths in \texttt{output/} are the Python QZ twin.
          \texttt{stoch\_simul} leaves the Taylor rule free and is
          \textbf{not} \textbf{(a) sterilised term premium}.
  \end{itemize}
\end{frame}

%----------------------------------------------------------------------
\begin{frame}{Three shocks, and how (a) is sterilised}
  \begin{center}
    \includegraphics[width=\textwidth,height=0.36\textheight,keepaspectratio]%
      {\irfdir bank_rate_sterilisation.pdf}
  \end{center}
  \vspace{-0.25em}
  {\scriptsize
  \textbf{(a) sterilised term premium:} anticipated Bank Rate peg,
  $H=12$ quarters
  (\texttt{H\_peg} in \texttt{src/calibration.py},
  \texttt{dynare/params\_calibration.m},
  \texttt{dynare/run\_uk\_long\_rate.m}).
  The Taylor rule stays in the model. For the first $H$ quarters
  $\varepsilon^{\mathrm{ster}}_t=-i^{TR}_t$, so}
  {\footnotesize
  \begin{align*}
    i_t
      &= \rho_i i_{t-1}
       + (1-\rho_i)\bigl(\varphi_\pi \pi_t + \tfrac14\varphi_y Y_t\bigr)
       + \psi_{\mathrm{nfa}} E_t nfa_{t+1}
       + \nu_t + \varepsilon^{\mathrm{ster}}_t
       = 0.
  \end{align*}
  }
  \vspace{-0.6em}
  {\scriptsize
  $\varphi_y=0.125$ is the annual convention, hence the factor $1/4$.
  Residuals are anticipated; the rule resumes after quarter~12.
  \textbf{(b) path news} is not sterilised.
  \textbf{(c) sovereign} is not yet modelled.}
\end{frame}

%----------------------------------------------------------------------
\begin{frame}{New-loan rate versus the mortgage stock}
  Debt service uses the stock rate. User cost and new borrowing use
  the new-loan rate.
  \begin{align*}
    R^{m}_{t}
      &= \omega_{S}\, i_{t} + \omega_{L}\, R^{L}_{t} + \zeta_{t},
      \qquad \omega_{S}+\omega_{L}=1, \\[0.15em]
    R^{m,\mathrm{stock}}_{t}
      &= (1-\lambda_{q})\, R^{m,\mathrm{stock}}_{t-1}
       + \lambda_{q}\, R^{m}_{t}.
  \end{align*}
  \begin{itemize}
    \item $\lambda_{q}=0.09$ is held fixed (sketch). Mean spell
          $1/0.09 \approx 11.1$ quarters; annual reset probability
          $1-0.91^{4}\approx 31\%$.
          Provisional weights: $\omega_{S}=0.30$, $\omega_{L}=0.70$.
    \item Borrower consumption:
          \[
            C^{b}_{t}
              = \alpha_{y} Y_{t}
              - \mu_{ds}\, R^{m,\mathrm{stock}}_{t}
              + \mu_{\mathrm{coll}}\lambda_{q}
                \bigl(P^{h}_{t}-B^{m}_{t-1}\bigr).
          \]
          $\mu_{ds}=15$, $\alpha_{y}=0.15$, $\mu_{\mathrm{coll}}=0.10$.
          The collateral term is provisional, and it only bites on the
          fraction $\lambda_{q}$ that remortgages.
  \end{itemize}
\end{frame}

%----------------------------------------------------------------------
\begin{frame}{Interest bill: index-linked share and duration}
  Do not double-count $s^{IL}$.
  $B^{\mathrm{nom}}=(1-s^{IL})B^{g}$ and $B^{IL}=s^{IL}B^{g}$,
  with $B^{g}$ equal to 4 quarters of GDP
  (about 100\% of annual GDP).
  In the linear units, $b^{\mathrm{nom}}=3$ and $b^{IL}=1$.
  \begin{align*}
    R^{\mathrm{eff}}_{t}
      &= (1-\delta_{D})\, R^{\mathrm{eff}}_{t-1} + \delta_{D}\, R^{L}_{t},
      \qquad \delta_{D}=\frac{1}{4D}, \\[0.2em]
    IB_{t}
      &= b^{\mathrm{nom}} R^{\mathrm{eff}}_{t}
       + b^{IL}\bigl(r^{IL}_{t}+\pi_{t}\bigr)
       + \frac{i_{ss}}{100}\widetilde{B}^{g}_{t}.
  \end{align*}
  \begin{itemize}
    \item $s^{IL}=0.25$ and $D=9.1$ years are held fixed, so
          $\delta_{D}=1/36.4\approx 0.02747$.
          On impact, from a zero lag,
          $R^{\mathrm{eff}}_{0}/R^{L}_{0}=1/(4D)$.
    \item The index-linked real coupon refixes toward $R^{L}-E\pi$.
          The inflation uplift sits \emph{inside} reported $IB$.
          A disinflation therefore cuts the bill on impact even while
          the conventional coupon is starting to refix.
  \end{itemize}
\end{frame}

%----------------------------------------------------------------------
\begin{frame}{UIP: a rise in $rer$ is a depreciation}
  \begin{equation*}
    rer_{t}
      = E_{t} rer_{t+1}
      - \bigl(i_{t}-E_{t}\pi_{t+1}\bigr)
      + \chi_{\mathrm{nfa}}\, nfa_{t}
      + \rho^{fx}_{t}.
  \end{equation*}
  \begin{itemize}
    \item $\chi_{\mathrm{nfa}}=0.01$. Net exports:
          $NX_{t}=\eta_{rer}\, rer_{t}-\eta_{y} Y_{t}+\eta_{y^{\ast}} Y^{\ast}_{t}$,
          with $\eta_{rer}=0.06$ and $\eta_{y}=0.18$.
    \item Under \textbf{(a) sterilised term premium}, $i_{t}=0$ on the peg
          while inflation falls, so the real short rate rises and
          sterling appreciates ($rer$ falls). That is the review's sign
          for (a). On this path $rer$ is $-1.875\%$ on impact
          (\texttt{irf\_tp\_sterilised.csv}).
    \item A passive tax rule $\tau_{t}=\varphi_{dg}\widetilde{B}^{g}_{t-1}$
          ($\varphi_{dg}=0.06$) closes the fiscal loop.
          $\varphi_{\pi}=1.5$, $\beta=0.995$.
  \end{itemize}
\end{frame}

%----------------------------------------------------------------------
\begin{frame}{Calibration: three cells held fixed, the rest provisional}
  {\footnotesize
  \setlength{\tabcolsep}{3.5pt}
  \renewcommand{\arraystretch}{1.02}
  \begin{tabular}{@{}lll>{\raggedright\arraybackslash}p{5.5cm}@{}}
    \toprule
    & Value & Status & Role \\
    \midrule
    $\lambda_{q}$ & 0.09 & \textbf{fixed} & quarterly remortgage hazard \\
    $s^{IL}$ & 0.25 & \textbf{fixed} & index-linked share of the gilt stock \\
    $D$ & 9.1y & \textbf{fixed} & modified duration \\
    $\kappa$ & 0.004 & \textbf{provisional} & NKPC slope. Flat, so a finite peg is not a Fisher spiral \\
    $\varphi_{h}$, $\kappa_{c}$ & 0.75, 0.055 & \textbf{provisional} & housing user cost; loading on $C^{b}$ \\
    $\kappa_{H}$, $\kappa_{\mathrm{level}}$ & 0.03, 0.05 & \textbf{provisional} & housing stock and level anchor \\
    $\mu_{\mathrm{coll}}$ & 0.10 & \textbf{provisional} & LTV / collateral, scaled by $\lambda_{q}$ \\
    $\omega_{S}$, $\omega_{L}$ & 0.30, 0.70 & provisional & Bank Rate and long-rate weights in $R^{m}$ \\
    $\beta$, $\varphi_{\pi}$, $\rho_{i}$ & 0.995, 1.5, 0.86 & provisional & discount factor and the Taylor rule \\
    $\varphi_{y}$, $\psi_{\mathrm{nfa}}$ & 0.125, 0.08 & provisional & annual output weight; NFA term inside $i^{TR}$ \\
    $\rho_{TP}$, $\rho_{\nu}$ & 0.90, 0.75 & provisional & persistence of (a) and of (b) \\
    $h_{s}$, $\sigma$, $\psi^{TP}$ & 0.40, 6, 2 & provisional & saver level rule; $\mathrm{eis}=(1-h_{s})/\sigma=0.10$ \\
    \bottomrule
  \end{tabular}
  }
  \vspace{0.15em}
  {\scriptsize
  Full list in \texttt{src/calibration.py} and the \texttt{.mod}.
  $\lambda_{q}$, $s^{IL}$ and $D$ are not free.
  $\kappa$ and the housing loadings are provisional: an earlier pass
  troughed near $-9.5\%$ per $+100$~bp. Not UK estimates.}
\end{frame}

%======================================================================
\section{Results}
%----------------------------------------------------------------------
\begin{frame}{(a) sterilised term premium}
  \begin{center}
    \includegraphics[width=\textwidth,height=0.62\textheight,keepaspectratio]%
      {\irfdir a_macro.pdf}
  \end{center}
  \vspace{-0.15em}
  {\scriptsize
  Output trough $-0.55\%$ in quarter~4.
  House prices trough $-2.6\%$ in quarter~8.
  Interest bill peaks at $+0.20$~pp of GDP in quarter~24.}
  \Calibrated
\end{frame}

%----------------------------------------------------------------------
\begin{frame}{(a) sterilised term premium: who cuts consumption}
  \begin{center}
    \includegraphics[width=\textwidth,height=0.38\textheight,keepaspectratio]%
      {\irfdir a_consumption_mortgage.pdf}
  \end{center}
  \vspace{-0.2em}
  \begin{itemize}
    \item Savers' consumption falls first: trough $-0.77\%$ in quarter~3.
          It does not rise later
          ($C^{s}$ stays in $[-0.77\%,-0.08\%]$).
    \item Borrowers' consumption falls more, and later: trough $-1.04\%$
          in quarter~9.
    \item Only about 9\% of the mortgage stock reprices each quarter
          ($\lambda_q=0.09$). New-loan $R^{m}$ jumps 70~bp on impact;
          the stock rate only 6.3~bp, and peaks in quarter~10,
          next to the borrower trough.
  \end{itemize}
  {\scriptsize Exact troughs in \texttt{success\_checks.txt}:
  $C^{s}$ $-0.768\%$ at q3, $C^{b}$ $-1.038\%$ at q9.}
  \Calibrated
\end{frame}

%----------------------------------------------------------------------
\begin{frame}{(a) sterilised term premium against (b) path news}
  \begin{center}
    \includegraphics[width=\textwidth,height=0.30\textheight,keepaspectratio]%
      {\irfdir a_vs_b.pdf}
  \end{center}
  \vspace{-0.15em}
  \begin{center}
  {\scriptsize
  \setlength{\tabcolsep}{5pt}
  \begin{tabular}{@{}lcc@{}}
    \toprule
    & (a) sterilised term premium & (b) path news \\
    \midrule
    $R^{L}$ on impact & $+100$~bp & $+100$~bp \\
    Bank Rate on impact & pegged at 0 ($H=12$) & $+315$~bp \\
    Output trough & $-0.55\%$ (q4) & $-1.178\%$ (q5) \\
    House-price trough & $-2.6\%$ (q8) & $-6.551\%$ (q7) \\
    \bottomrule
  \end{tabular}
  }
  \end{center}
  {\scriptsize
  (b), from \texttt{irf\_short\_rate\_news.csv}: weight $1/(4D)$ on the
  current short rate, so Bank Rate peaks at $+560$~bp in quarter~4.
  $R^{L}$ later dips to $-11.8$~bp in quarter~16 under (b) and stays
  positive under (a) (minimum $+2.35$~bp).}
  \Calibrated
\end{frame}

%----------------------------------------------------------------------
\begin{frame}{(a) sterilised term premium: interest-bill split}
  \begin{center}
    \includegraphics[width=\textwidth,height=0.32\textheight,keepaspectratio]%
      {\irfdir ib_decomposition.pdf}
  \end{center}
  \vspace{-0.35em}
  {\scriptsize
  Columns \texttt{ib\_nom} and \texttt{ib\_il} on
  \texttt{irf\_tp\_sterilised.csv}.
  $IB$ is percentage points of GDP.}
  \begin{itemize}
    \item Index-linked piece, including the uplift:
          $-0.0563$~pp of GDP on impact (disinflation).
    \item Nominal coupon: rises immediately, but only at rate
          $1/(4D)$. Peak $+0.1585$~pp in quarter~19.
          $R^{\mathrm{eff}}$ itself peaks in quarter~17.
    \item Sum: impact $-0.034$~pp, peak $+0.20$~pp of GDP in quarter~24
          (the check prints $+0.2017$).
          Coupon half-life is about 25 quarters; $\rho_{TP}=0.90$
          pulls the peak forward of that half-life.
  \end{itemize}
  \Calibrated
\end{frame}

%======================================================================
\section{Next steps}
%----------------------------------------------------------------------
\begin{frame}{Next: high-frequency gilt local projections}
  \texttt{src/lp\_proxyvar.py} is a scaffold. It does not ship a surprise file.
  With no \texttt{data/gilt\_surprises.csv} it writes
  \texttt{output/lp\_irf\_stub\_Y\_Ph.csv} full of NaNs.
  \begin{itemize}
    \item \textbf{(a) sterilised term premium,} empirical analogue:
          the high-frequency 10-year gilt change, residualised on a
          constant, the Bank Rate surprise, and the near-term OIS path
          (\texttt{orthogonalise\_gilt\_surprise}).
    \item \textbf{(b) path news:} the OIS path itself.
    \item Outcomes the stub is shaped for: ONS monthly GDP and a UK
          house-price index. Horizons 0--12 in the stub.
          Reference already in the code: Cesa-Bianchi, Thwaites and
          Vicondoa (2020).
    \item \textbf{(c) sovereign} is \textbf{not yet modelled}.
    \item \textbf{Next estimation step:} a linear Jord\`a baseline on
          that residualised surprise, then a state-dependent / non-linear
          variant \`a la David, Giacomini, Jiao and Wang.
          The non-linear variant is not in the repo yet.
  \end{itemize}
  {\scriptsize No UK surprise series is invented here.
  A synthetic DGP exists only inside \texttt{src/test\_lp\_algebra.py}.}
\end{frame}

%----------------------------------------------------------------------
\begin{frame}{Then: Dynare Bayesian estimation on UK data}
  \begin{itemize}
    \item Not run. The bottom of \texttt{uk\_long\_rate.mod} leaves
          \texttt{estimated\_params} for later.
          $\lambda_{q}$, $s^{IL}$ and $D$ stay fixed.
    \item A later block may free $\kappa$, $\rho_{TP}$, $\rho_{\nu}$,
          $\omega_{S}$, $\mu_{ds}$ and the housing loadings, and only
          once the observables and the high-frequency instruments
          are actually loaded.
    \item \textbf{Prior for $\kappa$.}
          Centre it on an import-penetration-dependent Phillips-curve
          slope \`a la Hottman and Reyes-Heroles, converted to an
          aggregate slope, and keep the prior \textbf{wide}.
          The calibrated $\kappa=0.004$ is a flat slope that stops a
          finite peg becoming a Fisher spiral. It is not that prior,
          and it is not a UK posterior.
    \item Sketch observable list, not yet loaded: Bank Rate / SONIA,
          the 10-year gilt and gilt--OIS, MLAR 2-year and 5-year fixes,
          the Bankstats stock rate, UK house prices, ONS consumption,
          investment, GDP and CPI, the sterling ERI, and the DMO
          interest bill.
  \end{itemize}
\end{frame}

%----------------------------------------------------------------------
\begin{frame}{Takeaways}
  \begin{itemize}
    \item Separate \textbf{(a) sterilised term premium},
          \textbf{(b) path news} and \textbf{(c) sovereign}.
          UK path-shock macro peaks are not UK term-premium IRFs.
          LDI / September 2022 stays a boxed non-linear stress effect.
    \item In the \textbf{calibrated} twin, a $+100$~bp
          \textbf{(a) sterilised term premium} ($H=12$) lowers output
          (trough $-0.55\%$ in quarter~4) and house prices
          (trough $-2.6\%$ in quarter~8). The interest bill peaks at
          $+0.20$~pp of GDP in quarter~24, after an index-linked dip.
          Savers' consumption falls first ($-0.77\%$ in quarter~3) and
          does not rise. Borrowers fall more, and later
          ($-1.04\%$ in quarter~9), because only about 9\% of the
          mortgage stock reprices each quarter ($\lambda_q=0.09$).
    \item \textbf{(b) path news}, scaled to the same long-rate impact,
          is a much larger short-rate path and a deeper contraction.
          \textbf{(c) sovereign} is not yet modelled.
    \item These are not estimates. $\kappa$ and the housing parameters
          are provisional. Next: residualised gilt local projections,
          including a state-dependent variant, then Bayesian estimation
          with a wide prior on $\kappa$.
  \end{itemize}
\end{frame}

%======================================================================
\section{Appendix}
%----------------------------------------------------------------------
\begin{frame}{Appendix: yield curve, saver rule, housing curve}
  {\footnotesize
  \begin{align*}
    R^{EH}_{t}
      &= \delta_{D}\, i_{t} + (1-\delta_{D})\, E_{t} R^{EH}_{t+1},
      \qquad
      R^{L}_{t} = R^{EH}_{t}+TP_{t}. \\[0.35em]
    C^{s}_{t}
      &= h_{s} C^{s}_{t-1}
       - \psi^{TP} TP_{t}
       - \mathrm{eis}\,(i_{t}-E_{t}\pi_{t+1})
       + \psi^{y} Y_{t}
       + \psi^{IB} IB_{t}
       - \psi^{\tau}\tau_{t}. \\[0.35em]
    (1+\beta+\kappa_{\mathrm{level}}) P^{h}_{t}
      &= P^{h}_{t-1}
       + \beta E_{t} P^{h}_{t+1}
       + \kappa_{c} C^{b}_{t} \\
      &\quad
       - \varphi_{h}\bigl(R^{m}_{t}-E_{t}\pi_{t+1}\bigr)
       - \kappa_{H} H_{t}.
  \end{align*}
  }
  \vspace{0.15em}
  {\scriptsize
  The saver equation is a level rule.
  A forward habit Euler plus a $\Delta TP$ capital-gain term had produced
  a positive tail in $C^{s}$ (peak about $+0.5\%$ on an earlier pass).
  $\psi^{TP}=2>0$ books the term premium as a duration loss.
  Housing loadings are provisional: an earlier pass troughed near
  $-9.5\%$ per $+100$~bp, which was far too large next to UK
  Bank Rate responses of a few percent.
  Equations are the \texttt{.mod} / \texttt{NOTES.md} system.}
\end{frame}

%----------------------------------------------------------------------
\begin{frame}{Appendix: peg length is not doing the work}
  \begin{center}
    \includegraphics[width=\textwidth,height=0.36\textheight,keepaspectratio]%
      {\irfdir peg_robustness.pdf}
  \end{center}
  \vspace{-0.35em}
  {\scriptsize
  \setlength{\tabcolsep}{3.5pt}
  \renewcommand{\arraystretch}{1.02}
  \begin{tabular}{@{}lcccccc@{}}
    \toprule
    $H$ & $Y$ trough & $P^{h}$ trough & $\max C^{s}$ & $C^{b}$ trough & $IB$ peak & $\min R^{L}$ \\
    \midrule
    8  & $-0.540\%$ q4 & $-2.540\%$ q8 & $-0.082\%$ & $-1.024\%$ & $+0.2016$ q25 & $+2.38$~bp \\
    12 & $-0.545\%$ q4 & $-2.597\%$ q8 & $-0.081\%$ & $-1.038\%$ & $+0.2017$ q24 & $+2.35$~bp \\
    20 & $-0.549\%$ q4 & $-2.637\%$ q9 & $-0.080\%$ & $-1.046\%$ & $+0.2036$ q24 & $+2.25$~bp \\
    40 & $-0.550\%$ q4 & $-2.645\%$ q9 & $-0.079\%$ & $-1.048\%$ & $+0.2046$ q24 & $+2.11$~bp \\
    \bottomrule
  \end{tabular}

  \vspace{0.2em}
  Only $H=12$ is \textbf{(a) sterilised term premium}.
  The figure labels the other pegs as robustness.
  Rounded from \texttt{output/irf\_peg\_robustness.csv}.
  Results slides round the $H=12$ row to output $-0.55\%$,
  house prices $-2.6\%$, and an interest bill of $+0.20$~pp.
  $\kappa=0.004$ and $\psi_{\mathrm{nfa}}$ are provisional.}
  \Calibrated
\end{frame}

%----------------------------------------------------------------------
\begin{frame}{Appendix: the unsterilised term premium is not (a)}
  Same $TP$ innovation as in the sterilised run, Taylor rule left on.
  This is what \texttt{stoch\_simul(e\_tp)} does.
  File: \texttt{output/irf\_tp\_unsterilised.csv}.
  \textbf{Not} \textbf{(a) sterilised term premium}.
  \vspace{0.4em}
  {\footnotesize
  \begin{tabular}{@{}ll@{}}
    \toprule
    $R^{L}$ on impact & $+94.9$~bp (not rescaled to 100) \\
    Bank Rate & $-6.0$~bp on impact; trough $-24.2$~bp in quarter~7 \\
    Output & impact $-0.349\%$; trough $-0.487\%$ in quarter~4 \\
    House prices & trough $-2.212\%$ in quarter~9 \\
    Saver consumption & stays negative (peak $-0.080\%$) \\
    \bottomrule
  \end{tabular}
  }
  \vspace{0.45em}
  {\scriptsize
  The endogenous easing is small next to \textbf{(b) path news},
  and the sign pattern for output, house prices and $C^{s}$ matches
  the sterilised path. It is still a different experiment:
  Bank Rate is not pegged.}
  \Calibrated
\end{frame}

%----------------------------------------------------------------------
\begin{frame}{Appendix: what the MVP leaves out}
  \textbf{In the linear baseline}
  \begin{itemize}
    \item \textbf{(a) sterilised term premium} with $H=12$,
          \textbf{(b) path news}, and the interest-bill block with
          $s^{IL}=0.25$.
    \item One mortgage hazard. Savers on a level rule. Flat provisional
          Phillips curve. Housing user-cost curve with provisional loadings.
  \end{itemize}
  \textbf{Not yet modelled}
  \begin{itemize}
    \item \textbf{(c) sovereign.}
  \end{itemize}
  \textbf{Not in the MVP}
  \begin{itemize}
    \item LDI / gilt-stress amplification. CRE, buy-to-let, a second
          mortgage cohort, and a corporate EBP block.
    \item Bayesian estimation, UK surprise data, and the state-dependent
          local projection.
  \end{itemize}
\end{frame}

%----------------------------------------------------------------------
\begin{frame}{Selected references}
  {\scriptsize
  Albuquerque, Lazarowicz and Lenney (2025), BoE SWP~1,115.\\
  Best et al.\ (2020), UK EIS, as cited in the review.\\
  Burr and Willems (2024), \emph{Bank of England Quarterly Bulletin}.\\
  Case, Quigley and Shiller (2001), NBER WP~8606.\\
  Cesa-Bianchi, Thwaites and Vicondoa (2020), \emph{European Economic Review}.\\
  Gertler and Karadi (2015), \emph{AEJ: Macroeconomics}.\\
  Gilchrist and Zakraj\v{s}ek (2007), NBER WP~13174;
  (2012), \emph{American Economic Review}.\\
  Gourinchas and Rey (2007), \emph{Journal of Political Economy}.\\
  Greenwood and Vayanos (2014), \emph{Review of Financial Studies}.\\
  Greenwood, Hanson, Stein and Sunderam (2023), \emph{Quarterly Journal of Economics}.\\
  Hall and Jorgenson (1967), \emph{American Economic Review}.\\
  Havranek et al.\ (2015), as cited in the review.\\
  IMF (2024), UK Selected Issues; IMF WEO April 2024, Ch.~2.\\
  Joyce and Lengyel (2024), BoE SWP~1,097.
  Kaminska, Mumtaz and \v{S}ustek (2021), BoE SWP~914.\\
  Mishkin (2007), Jackson Hole.
  Pinter (2023), BoE SWP~1,019.
  Rajan, Rodriguez-Tous and Salgado-Moreno (2025), BoE SWP~1,158.\\
  Anjum and Herler (2023), Joseph et al.\ (2019), McKay and Wieland,
  Shah et al.\ (2024), Slacalek (2009), Forbes et al.:
  \emph{cited via} Burr and Willems (2024).\\
  BoE Financial Policy Summary, 30~September 2022;
  Financial Stability Report, December 2022.\\
  Compile: \texttt{pdflatex uk\_long\_rate\_seminar.tex}
  from \texttt{uk-long-rate/beamer/}.
  }
\end{frame}

\end{document}
